\section*{Manipulador con indicador y cable de control}


\begin{center}
  \includegraphics[scale=0.5]{manipulador.png}
\end{center}

% \begin{figure}[h!]
%     \centering
%     \includegraphics[width=0.8\textwidth]{path/to/figure2.jpg} % Utilize la ruta adecuada para la imagen de la Figura 2.
%     \caption{Manipulador con indicador y cable de control}
%     \label{fig:manipulador}
% \end{figure}

% \begin{itemize}
%     \item[a)] Ensamble
%     \item[b)] Tapa final del conductor
%     \item[c)] Manija
%     \item[d)] Indicador
%     \item[e)] Conductores para cable
% \end{itemize}

Este conjunto (Figura 2) proporciona los medios para accionar la fuente desde el cabezal y a través del ducto, a la posición de exposición. Se cuenta con 15.2 metros de cable de acero inoxidable, los cuales son impulsados por un engranaje especial dentro de dos conductores exteriores. En un extremo del cable, el impulsor se encuentra un gancho el cual se conecta al cable de la fuente para su operación.

Girando la manija, gira el engranaje moviendo de esta manera al cable impulsor (y la fuente) fuera del cabezal. En la sección de la manija está montado un medidor, el cual proporciona una lectura que indica la distancia recorrida por la fuente a través del ducto. Este conjunto manipulador trasladará la fuente (fuera del cabezal) una distancia de 7 m.
